% Source-derived chapter 09: Applications
% Original source: pixtons-formula-abel-jacobi-picard-stack
\chapter{Applications}
\label{pixtons-formula-abel-jacobi-picard-stack-chapter-09}
\source{pixtons-formula-abel-jacobi-picard-stack}

\begin{remark}[Source section]
\label{pixtons-formula-abel-jacobi-picard-stack-chapter-09-source}
\source{pixtons-formula-abel-jacobi-picard-stack}
This chapter records the source section; no Lean declaration has yet been checked for this material.
\notready
\end{remark}

% The source section title was: Applications
\label{sec:applications}
\source{pixtons-formula-abel-jacobi-picard-stack}
\subsection{Proofs of Theorem \ref{FPA1} and Conjecture A} \label{Sect:ProofConjA}
\source{pixtons-formula-abel-jacobi-picard-stack}
We start by recalling notions presented in Section \ref{Sect:Twistdiffintro}, but now in the more general setting of $k$-differentials.
Let $A=(a_1,\ldots,a_n)$ be a vector of zero and pole multiplicities
satisfying
$$\sum_{i=1}^n a_i= k(2g-2)\, .$$
Let $\HH^k_g(A) \subset \mathcal{M}_{g,n}$ be the closed (generally non-proper) locus
of pointed curves $(C,p_1,\ldots,p_n)$ satisfying the condition
$$\mathcal{O}_C\Big(\sum_{i=1}^n a_i p_i\Big) \simeq \omega_C^{\otimes k}\, .$$
In other words, $\HH^k_g(A)$ is the locus of (possibly) meromorphic $k$-differentials
with zero and pole multiplicities prescribed by $A$.
In \cite{Farkas2016The-moduli-spac}, a compact moduli space of twisted $k$-canonical divisors
$$\widetilde{\HH}^k_g(A) \subset \overline{\mathcal{M}}_{g,n}$$
is constructed extending $\HH^k_g(A) = \widetilde{\HH}^k_g(A) \cap \mathcal{M}_{g,n} $ to the boundary of $\overline{\mathcal{M}}_{g,n}$.

For $k \geq 1$ and $A$ not of the form $A = k \cdot A'$ with a vector $A'$ of nonnegative integers, the locus
$\widetilde{\HH}^k_g(A)$ is of pure codimension $g$ in
$\overline{\mathcal{M}}_{g,n}$ by \cite[Theorem 3]{Farkas2016The-moduli-spac} (for $k=1$) and \cite[Theorem 1.1]{Schmitt2016Dimension-theor} (for $k>1$).
A weighted fundamental cycle of $\widetilde{\HH}^k_g(A)$,
\begin{equation}\label{wwww4k}
\source{pixtons-formula-abel-jacobi-picard-stack}
\mathsf{H}^k_{g,A}\in \mathsf{CH}_{2g-3+n}(\oM_{g,n})\ ,
\end{equation}
is constructed in
\cite[Appendix A]{Farkas2016The-moduli-spac} and \cite[Section 3.1]{Schmitt2016Dimension-theor} with explicit nontrivial
weights on the irreducible components.
The closure
$$\overline{\HH}^k_{g,A}\subset \oM_{g,n}$$
contributes to the weighted fundamental class $\mathsf{H}^k_{g,A}$ with
multiplicity 1, but there are additional boundary contributions, as described in the references above.

The weighted fundamental class
$\mathsf{H}^k_{g,A}$
was conjectured in \cite{Farkas2016The-moduli-spac,Schmitt2016Dimension-theor}
to equal  the class given by Pixton's formula for the double ramification cycle. To state the conjecture, consider the shifted\footnote{The shift is needed since Pixton's original formula worked with powers of the log-canonical line bundle $\omega_C^{\text{log}} = \omega_C(\sum_{i=1}^n p_i)$ instead of $\omega_C$.} vector $\widetilde A = (a_1 + k, \ldots, a_n +k)$.

\begin{conjectureA}
For $k \geq 1$ and $A$ not of the form $A = k \cdot A'$ with a vector $A'$ of nonnegative integers, we have an equality
\[\mathsf{H}^k_{g,A} = 2^{-g} P_g^{g,k}(\widetilde A)\, ,\]
where $P_g^{g,k}(\widetilde A)$ is Pixton's cycle class defined in \cite[Section 1.1]{Janda2016Double-ramifica}.
\end{conjectureA}

 By combining Theorem \ref{Thm:main}  with previous results of \cite{Holmes2019Infinitesimal-s}, we can now prove the conjecture.

\begin{theorem}
\source{pixtons-formula-abel-jacobi-picard-stack}
\notready\label{FPA}
Conjecture A is true.
\end{theorem}
\begin{proof}
By Theorem 1.1 of \cite{Holmes2019Infinitesimal-s}, the weighted fundamental class $\mathsf{H}^k_{g,A}$ is equal to the double ramification cycle $\mathsf{DR}_{g,A, \omega^k}$ constructed in \cite{Holmes2017Extending-the-d}. By Theorem \ref{cccc},
$\mathsf{DR}_{g,A, \omega^k}$
is in turn given by the action of $\mathsf{DR}^{\mathsf{op}}_{g,A}$ on the fundamental class
of $\overline{\mathcal{M}}_{g,n}$ via the morphism $\phi_{\omega_\pi^k}: \overline{\mathcal{M}}_{g,n} \to \Picabs_{g,n,k(2g-2)}$ associated to the family
\begin{equation*}
\pi: \mathcal{C}_{g,n} \rightarrow \overline{\cM}_{g,n}\, , \ \ \ \ \omega_\pi^{k} \rightarrow \mathcal{C}_{g,n}\, .
\end{equation*}
By Theorem \ref{Thm:main}, the class $\mathsf{DR}^{\mathsf{op}}_{g,A}$ is computed by the tautological class
$$\P_{g,A,d}^g \in {\mathsf{CH}}^g_{\mathsf{op}}(\Picabs_{g,n,k(2g-2)})\, . $$ By Proposition \ref{Prop:PixtonPullbackCompatibility}, the action of $\P_{g,A,d}^g$ on $[\overline{\mathcal{M}}_{g,n}]$ is indeed given by Pixton's original formula $ 2^{-g} P_g^{g,k}(\widetilde A)$, finishing the proof. \end{proof}

The steps of the proof of Theorem \ref{FPA} are summarized as follows:
\begin{align*}
    \mathsf{H}^k_{g,A}&=\mathsf{DR}_{g,A, \omega^k} &&\cite[\text{Theorem }1.1]{Holmes2019Infinitesimal-s}\\
    &=\mathsf{DR}^{\mathsf{op}}_{g,A}(\phi_{\omega_\pi^k})([\overline{\mathcal{M}}_{g,n}]) &&\text{Theorem }\ref{cccc}\\
    &=\P_{g,A,d}^g(\phi_{\omega_\pi^k})([\overline{\mathcal{M}}_{g,n}]) &&\text{Theorem }\ref{Thm:main}\\
    &=2^{-g} P_g^{g,k}(\tilde A) &&\text{Proposition }\ref{Prop:PixtonPullbackCompatibility}\, .
\end{align*}
The result provides a completely geometric representative
of Pixton's cycle in terms of twisted $k$-differentials.
Theorem \ref{FPA1} of Section \ref{Sect:Twistdiffintro} is the $k=1$ case of
Theorem \ref{FPA}.

% A different
% approach to Conjecture A will be pursued in [Felix et al] \jocomment{To fill in.}
% using new logarithmic compactifications of the moduli of
% meromorphic $k$-differentials and the virtual localization formula
% in the log context.

%Theorem \ref{FPA} is proven in Section \ref{Sect:ProofConjA} where the parallel conjectures \cite{Schmitt2016Dimension-theor}
%for higher differentials  are also proven (by parallel  arguments).

\subsection{Closures}\label{clozzz}
\source{pixtons-formula-abel-jacobi-picard-stack}
Let $A=(a_1,\ldots, a_n)$ be a vector of integers satisfying
$$\sum_{i=1}^n a_i =k(2g-2)\, .$$
A careful investigation of the closure%\{ the superscripts escape from the $[-]$ a bit. We could use $\left[ ... \right]$ instead, but would maybe just make things ugly. } \jocomment{It's true that the superscript is not under the bar, but for me that would be ok (and it is the notation I had in my paper).}Fine by me, removing comments.
$$\HH^k_g(A) \subset \overline{\HH}^k_g(A) \subset \oM_{g,n}$$
is carried out in \cite{Bainbridge2018Compactificatio,Bainbridge2019Strata-of-k-dif}.
By a simple method presented in \cite[Appendix A]{Farkas2016The-moduli-spac} and \cite[Section 3.4]{Schmitt2016Dimension-theor},
Theorem \ref{FPA} easily determines the cycle classes of the
closures
$$ [\overline{\HH}^k_g(A)] \in \mathsf{CH}_{*} (\oM_{g,n})\, .$$
for the cases
\begin{itemize}
    \item $k=1$ and all $a_i$ nonnegative, when $\overline{\HH}^k_g(A)$ has pure codimension $g-1$ and
    \item $k\geq 1$ and $A$ not of the form $A = k \cdot A'$ with a vector $A'$ of nonnegative integers, when $\overline{\HH}^k_g(A)$ has pure codimension $g$. %\{Above we write the condition as `not a $k$-multiple of something non-negative'. Better to be uniform with this?}
\end{itemize}
In particular, from the recursive formula for $[\overline{\HH}^k_g(A)]$ and the fact that Pixton's cycle on $\oM_{g,n}$ is tautological, the following is immediate.
\begin{corollary}
\source{pixtons-formula-abel-jacobi-picard-stack}
\notready\label{kk333}
The cycles $[\overline{\HH}^k_g(A)]$ are tautological classes in $\mathsf{CH}_*(\oM_{g,n})$.
\end{corollary}
In the case $k=1$, Corollary \ref{kk333} was known by work of Sauvaget  \cite{Sauvaget2017Cohomology-clas}, who gave a different approach to $[\overline{\HH}^1_g(A)]$ in terms of tautological classes.
The recursive formulas for $[\overline{\HH}^k_g(A)]$ from
Corollary \ref{kk333} have been implemented\footnote{In the ongoing project \cite{CMZ20}, the authors study formulas for Euler characteristics of strata of differentials in terms of intersection numbers on the compactification of these strata constructed in \cite{BCGGM3}. The
implementation of $[\overline{\HH}^1_g(A)]$ has played
a role in corroborating their formulas.}
in the software \cite{admcycles} for computations in the tautological ring of $\Mbar_{g,n}$.

Another application of Conjecture A is presented in the recent paper \cite{sauvaget_volumes} by Sauvaget. The paper studies moduli spaces of flat surfaces of genus $g$ with conical singularities at marked points $p_1, ..., p_n$. The singularities have fixed cone angles $2 \pi \alpha_i$, for $\alpha_1, \ldots, \alpha_n \in \mathbb{R}$, summing to $2g-2+n$.
If all $\alpha_i$ are rational,
the spaces of flat surfaces naturally contain $\mathcal{H}_g^k(k A )$,
for  $$A=(\alpha_i-1)_{i=1}^n\, ,$$
as closed subsets (for $k$ sufficiently divisibly). These subsets equidistribute (with respect to natural measures) as $k \to \infty$. Using the
equidistribution,
Sauvaget is able to apply the recursive expression for $\overline{\HH}_g^k(k A)$ from Conjecture A to derive an explicit recursion for the volumes of the moduli spaces of flat surfaces.

\subsection{\texorpdfstring{$k$}{k}-twisted DR cycles with targets} \label{Sect:ktwistedDR}
\source{pixtons-formula-abel-jacobi-picard-stack}
We define  here  {\em $k$-twisted double ramification cycles with targets} via the class $\DRop_{g,A}$.

Let $X$ be a nonsingular projective variety with line bundle $\ca L$ and
an effective curve class
$\beta \in H_2(X, \mathbb{Z})$. Let
$$d_\beta= \int_\beta c_1(\ca L)\, .$$
Let $k\in \mathbb{Z}$ and $A=(a_1,\ldots, a_n)\in \mathbb{Z}^n$ satisfy
\begin{equation*}
d_\beta + k(2g-2+n) = \sum_{i=1}^{n}a_i\, .
\end{equation*}
Consider the morphism
\[
\phi\colon \MbarX \to \Picabs_{g,n,d_\beta+k(2g-2+n)}
\]
defined by the universal data
\begin{equation}\label{fftt666774}
\source{pixtons-formula-abel-jacobi-picard-stack}
\pi: \mathcal{C}_{g,n,\beta} \rightarrow \overline{\cM}_{g,n}(X,\beta)\, , \ \ \ \
f^*\ca L\otimes \omega^{\otimes k}_{log}
\rightarrow \mathcal{C}_{g,n,\beta}\, ,
\end{equation}
where
$f: \mathcal{C}_{g,n,\beta} \to X$
is the universal map.

%\[
%(f:(C,p_1,\cdots,p_n)\to X) \mapsto (C, p_1, \ldots, p_n, f^{*}\ca L \otimes %\omega^{\otimes k}_{log}(-\sum_{i=1}^{n}a_i p_i)).
%\]
\begin{definition}
\source{pixtons-formula-abel-jacobi-picard-stack}
\notready
The \textit{$k$-twisted $X$-valued double ramification cycle} is defined by
\[
\DR_{g,n,\beta}^{k}(X,\ca L)=\mathsf{DR}^{\mathsf{op}}_{g,A} (\phi) ( [\MbarX]^\textup{vir}) \in \Chow_{\vdim(g,n,\beta)-g}(\MbarX)\,.
\]
\end{definition}
In the notation of \cite[Section 0.4]{Janda2018Double-ramifica}, let $\P^{c,k,r}_{g,A,\beta}(X,\ca L)$ be the codimension $c$ part of the following expression
\begin{multline*}
\sum_{
\substack{\Gamma\in \G_{g,n,\beta}(X) \\
w\in \mathsf{W}_{\Gamma,r,k}(X)}
}
\frac{r^{-h^1(\Gamma)}}{|\Aut(\Gamma)| }
\;
j_{\Gamma_\delta*}\Bigg[
\prod_{i=1}^n \exp\left(\frac12 a_i^2 \psi_i + a_i \xi_i \right)\\
\prod_{v \in \V(\Gamma)} \exp\left(-\frac12 \eta(v) - k \eta_{1,1}(v) - \frac{k^2}{2} \eta_{0,2}(v)    \right)
\\
\prod_{e=(h,h')\in \E(\Gamma)}
\frac{1-\exp\left(-\frac{w(h)w(h')}2(\psi_h+\psi_{h'})\right)}{\psi_h + \psi_{h'}} \Bigg]\,.
\end{multline*}
The definition of the admissible $k$-weightings $w\in \mathsf{W}_{\Gamma,r,k}(X)$ is similar to that in \cref{Ssec:weightings} but with the condition (iii) replaced by
\[
k(2g(v)-2+n(v)) + \int_{\beta(v)} c_1(\ca L) = \sum_{v(h)=v}w(h)\,\text{ for }v \in \V(\Gamma)\, .
\]
As in the case $k=0$ discussed in \cite[Proposition 1]{Janda2018Double-ramifica}, the class $\P^{c,k,r}_{g,A,\beta}(X,\ca L)$ is polynomial in $r$ for all sufficiently large $r$.
Denote by $\P^{c,k}_{g,A,\beta}(X,\ca L)$ the value at $r=0$ of this polynomial.
%\{Do we need to justify that it is polynomial? Or cite the mysterious Pixton paper? }\jocomment{I added a reference to the $k=0$ version of this result. Actually, this polynomiality (in $r$) is really proven in \cite[Appendix]{Janda2016Double-ramifica}, the mystery polynomiality is the one in the $a_i$.} Thanks! Deleting comments.

By \cref{Thm:main} and a slight generalization of the procedure for pulling back Chow cohomology classes from $\Picabs_{g,n,d_\beta+k(2g-2+n)}$ to $\MbarX$ described in \cite[Section 1.5]{Janda2018Double-ramifica}, we have
\begin{eqnarray*}
    \DR_{g,n,\beta}^{k}(X,\ca L)&= &\mathsf{DR}^{\mathsf{op}}_{g,A}(\phi)( [\MbarX]^\textup{vir})\\
    &=&\P_{g,A,d_\beta+k(2g-2+n)}^g (\phi)( [\MbarX]^\textup{vir})\\
    &=&\P^{g,k}_{g,A,\beta}(X,\ca L)\, .
\end{eqnarray*}

\subsection{Proof of Theorem \ref{Thm:mainvan}}
For all $c>g$, we will prove
\begin{equation} \label{eqn:PgAdc_vanishing}
\source{pixtons-formula-abel-jacobi-picard-stack}
    \P_{g,A,d}^c\ = 0 \
\in {\mathsf{CH}}^c_{\mathsf{op}}(\Picabs_{g,n,d})\, .
\end{equation}
The path is parallel to the proof of Theorem \ref{Thm:main}.

%The proof only use the invariants of the expression $\P_{g,A,d}^c$ in \cref{sec:proof_of_main_theorem}.
By definition, the claim is
equivalent to showing that the map
\begin{equation}\label{ff55ff3}
\source{pixtons-formula-abel-jacobi-picard-stack}
 \P_{g,A,d}^c(\phi) : \Chow_*(B) \to \Chow_{*-c}(B)
 \end{equation}
is zero for every
morphism $\phi:B \to \Picabs_{g,n,d}$ from an (irreducible) finite type scheme $B$
corresponding to the data
$$C \to B\, , \ \ \  \ca L \to C\,.$$
Retracing the steps of \cref{sec:proof_of_main_theorem} (and using the invariance \cref{lem:N_invariance_P}  for the codimension $c$ part $\P_{g,A,d}^c$ of Pixton's formula), we can reduce to the situation where $\ca L$ on $C$ is relatively sufficiently
positive with respect to $C \to B$. As in \cref{sec:positive_line_bundle},
we can then find $$\psi: U_l \to B$$
such that $\psi^*$ is injective on Chow groups and such that the composition
$$U_l \to B \to \Picabs_{g,n,d}$$
factors through $\overline{\mathcal{M}}_{g,n}(\PP^l, d)'$. By Theorem 3.2
of \cite{Bae},
we have the vanishing $$\P_{g,A,d}^c (\PP^l,\ca O(1)) = 0
\, \in \, \mathsf{CH}_{\textup{vdim}(g,n,d)-c}(\overline{\mathcal{M}}_{g,n}(\PP^l, d))\, .$$
The same combination of \cref{lem:op_eq_ch_for_smooth} and the injectivity of $\psi^*$ then shows the desired vanishing of the map \eqref{ff55ff3}. \qed

% As in \cref{sec:proof_of_main_theorem} we can find an alteration $B' \to B$ together with a partial stabilization $f: C' \to C \times_B B'$ such that $C'$ admits disjoint sections meeting every component of every geometric fibre over $B'$. Applying \cref{lem:N_invariance_P} it suffices to consider the family $C' \to B'$ and the line bundle $f^* \ca L$. Here, by the arguments in \cref{sec:many_markings} we

%For $X=\PP^k$, $\ca L = \ca O(1)$ and $c>g$, we have vanishing $\P_{g,A,d}^c (X,\ca L) = 0$   by \cite{Bae}. The arguments of \cref{sec:proof_of_main_theorem} then imply


\subsection{Connections to past and future results}\label{pasfut}
\source{pixtons-formula-abel-jacobi-picard-stack}
The relations of Theorem \ref{Thm:mainvan}
 generalize several previous results.
For $g=0$ and $c=1$, the vanishing \eqref{eqn:PgAdc_vanishing}
was observed in \cite[Proposition 1.2]{deJongStarr}.
In fact, in genus $0$, there are many connections to past equations, see \cite[Section 4]{Bae}
for a full discussion with many examples including classical
equations and
the relations of \cite{LeePand}.

%In fact, all genus zero $\mathsf{DR}$ relations can be obtained from tautological relations on $\mathfrak{M}_{0,n}$.

Randal-Williams \cite{RandalWilliams2012}  proves
a vanishing result in  cohomology with integral coefficients on the locus $\Picabs_{g,0,d}^\textup{sm}$ of smooth curves
for every $d \in \mathbb{Z}$.
We can recover a version of Randal-William's
vanishing in operational Chow  with $\mathbb{Q}$-coefficients which extends to all of $\Picabs_{g,0,d}$.  By \cref{Lem:Pixformulafactorization} and \cref{Lem:vertterm},  Pixton formula's on the locus $\Picabs_{g,0,0}^\textup{sm}$ takes the simple form
\[\P_{g,\emptyset,0}^c = \frac{1}{c!} (\P_{g,\emptyset,0}^1)^c\, ,\ \ \  \P_{g,\emptyset,0}^1 = - \frac{1}{2} \pi_* (c_1(\ca L)^2)\]
for the universal curve and the universal line bundle
$$\pi: \frak C \to \Picabs_{g,0,0}\, , \ \ \ \ca L \to \frak C\, .$$
We claim, up to scaling, the relation
$$\Omega^{g+1}=0$$ of
\cite[Theorem A]{RandalWilliams2012} is exactly the restriction of the pullback of the relation $$(\P_{g,\emptyset,0}^1)^{g+1} = (g+1)! \P_{g,\emptyset,0}^{g+1} = 0
\in {\mathsf{CH}}^{g+1}_{\mathsf{op}}(\Picabs_{g,0,0})$$
under the morphism
\[\Picabs_{g,0,d} \to \Picabs_{g,0,0}\, ,\ \ \ (C, \ca L) \mapsto (C, \ca L^{\otimes 2g-2} \otimes \omega_C^{\otimes (-d)})\, .\]
Indeed, over the locus of smooth curves, the pullback of $\P_{g,\emptyset,0}^1$ is given by
\begin{multline*}
- \frac{1}{2} \pi_* (c_1(\ca L^{\otimes 2g-2} \otimes \omega_C^{\otimes (-d)})^2)=\\ -\frac{1}{2} \left((2g-2)^2 \pi_*(c_1(\ca L)^2) - 2d(2g-2)\pi_*(c_1(\ca L ) c_1(\omega_\pi)) + d^2 \pi_*(c_1(\omega_\pi)^2) \right)\, ,\end{multline*}
which matches the definition of $\Omega$ given in \cite[Theorem A]{RandalWilliams2012}
up to scalars.
% When we restrict our relations to the locus where the curve is smooth, these relations for higher genus cases were first observed in topological context, see \cite{RandalWilliams2012}.

In Gromov-Witten theory,  pulling back \eqref{eqn:PgAdc_vanishing} under the morphisms $$\MbarX \to \Picabs_{g,n,d}$$ described in  \cref{Sect:ktwistedDR} and capping with the virtual class $[\MbarX]^\textup{vir}$ simply recovers the known vanishing
\begin{equation} \label{eqn:PXbetavanishing}
\source{pixtons-formula-abel-jacobi-picard-stack}
    \P_{g,A,\beta}^{c,k}(X,\ca L) = 0 \in \Chow_{\vdim(g,n,\beta)-c}(\MbarX)
\end{equation}
for $c>g$ proven in \cite{Bae}. However, there are new applications for \emph{reduced} Gromov-Witten theory. Indeed, for a target $X$ having a nondegenerate holomorphic $2$-form, the virtual class of $\MbarX$ vanishes when $\beta \neq 0$.
To define invariants for such targets, the reduced class
\[[\MbarX]^\textup{red} \in \Chow_{*}(\MbarX)\]
is used instead, see \cite{bryanleung,maulikpandharipande}. By
pulling back \eqref{eqn:PgAdc_vanishing} and capping with $[\MbarX]^\textup{red}$, we obtain new relations among reduced Gromov-Witten invariants. An application to the Gromov-Witten theory of K3 surfaces will appear in \cite{BaeBuelles} related to
conjectures of \cite{ObPand}.


% This result has an application for reduced Gromov-Witten theory of $X$ which is not immediately obvious from previous works \cite{Bae,cj}.

% When $X$ has a nondegenerate holomorphic $2$ form, the virtual class of $\MbarX$ vanishes when $\beta \neq 0$, but the reduced class may not vanish. If we pullback $\mathsf{DR}$ relations to $\MbarX$ and acts on the reduced class, we get new relations among reduced Gromov-Witten invariants. An application to Gromov-Witten theory of K3 surfaces will appear in \cite{BaeBuelles}.


 %\{Looks very nice! On the relations of Randall-Williams, maybe worth mentioning that our results are stronger as we prove them in Chow, but weaker because we work with rational coefficients, so we cannot given information on the annihilator?}

%\textcolor{brown}{Y: In Rahul's note\\ \url{https://people.math.ethz.ch/~rahul/MRW.pdf}\\ he said the statement for rational Chow group was proven by Van der Geer. So, our result is stronger than Randal-Williams' in the sense that we proved the result on $\Picabs$. I agree that his result is sensitive in torsion part.}

\def\cprime{$'$}


 {Department of Mathematics, ETH Z\"urich} \\
 {younghan.bae@math.ethz.ch}\\

 {Mathematisch Instituut, Universiteit Leiden} \\
 {holmesdst@math.leidenuniv.nl}\\

 {Department of Mathematics, ETH Z\"urich} \\
 {rahul.pandharipande@math.ethz.ch}\\

 {Mathematical Institute, University of Bonn} \\
 {schmitt@math.uni-bonn.de}\\

 {Mathematisch Instituut, Universiteit Leiden} \\
 {r.m.schwarz@math.leidenuniv.nl}\\
